\documentclass[aspectratio=169]{beamer}
\usepackage{etex} % fixes new-dimension error
\usepackage{lmodern}
\usepackage[T1]{fontenc}

\usepackage{graphicx,amsmath}
\usepackage{stmaryrd} % cf. interleave
\usepackage{booktabs}
\usepackage{amscd}
\usepackage{multicol}
\usepackage[absolute,overlay]{textpos}
\usepackage{alltt}
\usepackage{proof}
\usepackage{subcaption}
\usepackage{tikz}
\usepackage{tikz-cd}
\usepackage[new]{old-arrows}
\usepackage[all]{xy}
\usepackage{pgfplots}
\usepackage{textcomp}
\usepackage{cancel}

\usepackage{transparent}
\usepackage{xspace}
\usepackage{listings}
\usepackage{pdfpages}
\usepackage{relsize}



%%%%%%%%%%%%% Macros
% \newcommand{\Ban}{\catfont{Ban}}
% \newcommand{\Cats}{\catfont{Cat}}

%%%% Misc
%% Operations
\newcommand{\sem}[1]{\llbracket #1 \rrbracket}
\DeclareMathOperator{\img}{\mathrm{im}}
\DeclareMathOperator{\dom}{\mathrm{dom}}
\DeclareMathOperator{\codom}{\mathrm{codom}}
%% Sets of numbers
\newcommand{\N}{\mathbb{N}}
\newcommand{\Z}{\mathbb{Z}}
\newcommand{\Nats}{\mathbb{N}}
\newcommand{\Reals}{\mathbb{R}}
\newcommand{\Rz}{\Reals_{\geq 0}}
\newcommand{\Complex}{\mathbb{C}}
%% Writing
\newcommand{\cf}{\emph{cf.}}
\newcommand{\ie}{\emph{i.e.}}
\newcommand{\eg}{\emph{e.g.}}
\newcommand{\df}[1]{\emph{\textbf{#1}}}
%%%%%%%%%%%%%%%% End of Misc

%%%% Programming Stuff
%% Types
\newcommand{\typefont}[1]{\mathbb{#1}}
\newcommand{\typeOne}{1}
\newcommand{\typeTwo}{2}
\newcommand{\typeA}{\typefont{A}}
\newcommand{\typeX}{\typefont{X}}
\newcommand{\typeB}{\typefont{B}}
\newcommand{\typeC}{\typefont{C}}
\newcommand{\typeV}{\typefont{V}}
\newcommand{\typeD}{\typefont{D}}
\newcommand{\typeI}{\typefont{I}}
%% RuleName
\newcommand{\rulename}[1]{(\mathrm{#1})}
%% Sequents
\newcommand{\jud}{\vdash}
\newcommand{\vljud}{\rhd}
\newcommand{\cojud}{\vdash_{\co}}
\newcommand{\vl}{\mathtt{v}}
\newcommand{\co}{\mathtt{c}}
% Program font
\newcommand{\prog}[1]{\mathtt{#1}}
\newcommand{\pseq}[3]{#1 \leftarrow #2; #3}
\newcommand{\ppm}[4]{(#1,#2) \leftarrow #3; #4}
\newcommand{\pinl}[1]{\prog{inl}(#1)}
\newcommand{\pinr}[1]{\prog{inr}(#1)}
\newcommand{\pcase}[4]{\prog{ case } #1 \prog{ of } \pinl{#2} \Rightarrow #3 ; \pinr{#2} \Rightarrow #4}
%% Sets of terms
\newcommand{\ValuesBP}[2]{\mathsf{Values}(#1, #2)}
\newcommand{\TermsBP}[2]{\mathsf{Terms}(#1, #2)}
\newcommand{\closValP}[1]{\ValuesBP{\emptyset}{#1}}
\newcommand{\closTermP}[1]{\TermsBP{\emptyset}{#1}}
\newcommand{\closVal}{\closValP{\typeA}}
\newcommand{\closTerm}{\closTermP{\typeA}}
%% Contextual equivalence
\newcommand{\ctxeq}{\equiv_{\prog{ctx}}}
%%%% End of Programming Stuff


% Misc by José
\newcommand{\wrap}[2][]{\begin{tabular}[#1]{@{}c@{}}#2\end{tabular}}
\newcommand{\mwrap}[1]{\ensuremath{\begin{array}{@{}c@{}}#1\end{array}}}
\def\trans#1{\xrightarrow{#1}}  % - a - > 
\def\Trans#1{\stackrel{#1}{\Longrightarrow}} % =a=> 
\newcommand{\transp}[2][35]{\color{fg!#1}#2}
\newcommand{\transpt}[2][.35]{\tikz{\node[inner sep=1pt,fill opacity=0.5]{#2}}}
\newcommand{\faded}[2][0.4]{{\transparent{#1}#2}} % alternative to "transp" using transparent package
\newcommand{\set}[1]{\left\{ #1 \right\}} % {a,b,...z}
\newcommand{\mi}[1]{\ensuremath{\mathit{#1}}\xspace}
\newcommand{\mf}[1]{\ensuremath{\mathsf{#1}}\xspace}
% \newcommand{\gold}[1]{\textcolor{darkgoldenrod}{#1}\xspace}


%------ using color ---------------------------------------------------------
\definecolor{goldenrod}{rgb}{.80392 .60784 .11373}
\definecolor{darkgoldenrod}{rgb}{.5451 .39608 .03137}
\definecolor{brown}{rgb}{.15 .15 .15}
\definecolor{darkolivegreen}{rgb}{.33333 .41961 .18431}
\definecolor{myGray}{gray}{0.85}
%
%
\newcommand{\red}[1]{\textcolor{red!80!black}{#1}\xspace}
\newcommand{\blue}[1]{\textcolor{blue}{#1}\xspace}
\newcommand{\gold}[1]{\textcolor{darkgoldenrod}{#1}\xspace}
\newcommand{\gray}[1]{\textcolor{myGray}{#1}\xspace}
% \def\alert#1{{\darkgoldenrod #1}}
% \def\alert#1{{\alert{#1}}}
%\def\brw#1{{\brown #1}}
% \def\structure#1{{\blue #1}}
% \def\tstructure#1{\textbf{\darkblue #1}}
%%\def\gre#1{{\green #1}}
\def\gre#1{{\darkolivegreen #1}}
\def\gry#1{{\textcolor{gray}{#1}}}
\def\rdb#1{{\red #1}}
\def\st{\mathbf{.}\,}
\def\laplace#1#2{*\txt{\mbox{ \fcolorbox{black}{myGray}{$\begin{array}{c}\mbox{#1}\\\\#2\\\\\end{array}$} }}}
%\newcommand{\galois}[2]{#1\; \dashv\; #2}



\def\ainv#1{\overline{#1}}
\def\aconv#1{#1^{\circ}} 
\def\rtran#1{\stackrel{#1}{\longrightarrow}}
\def\pair#1{\langle #1 \rangle}
\def\setdef#1#2{\mathopen{\{} #1 \asor #2 \mathclose{\}}}
\def\imp{\mathbin{\Rightarrow}}
\def\dimp{\mathbin{\Leftrightarrow}}
\def\rimp{\mathbin{\Leftarrow}}
\def\rra{\longrightarrow}
% \def\rcb#1#2#3#4{\def\nothing{}\def\range{#3}\mathopen{\langle}#1 \ #2 \ \ifx\range\nothing::\else: \ #3 :\fi \ #4\mathclose{\rangle}}
\def\abv{\stackrel{\rm abv}{=}}



% Spliting frames in 2 columns
\newcommand{\splittwo}[4]{ 
  \begin{columns}[T]% align columns
  \begin{column}{#1\textwidth} #3 \end{column} ~~~
  \begin{column}{#2\textwidth} #4 \end{column} \end{columns}
}\newcommand{\frsplit}[3][.48]{
  \begin{columns}%[T] % align columns
  \begin{column}{#1\textwidth} #2 \end{column} ~~~
  \begin{column}{#1\textwidth} #3 \end{column} \end{columns}
}
\newcommand{\frsplitdiff}[5][]{
  \begin{columns}[#1]%[T] % align columns
  \begin{column}{#2\textwidth} #4 \end{column} ~~~
  \begin{column}{#3\textwidth} #5 \end{column} \end{columns}
}
\newcommand{\frsplitt}[3][.48]{
  \begin{columns}[T] % align columns
  \begin{column}{#1\textwidth} #2 \end{column} ~~~
  \begin{column}{#1\textwidth} #3 \end{column} \end{columns}
}
\newcommand{\col}[2][.48]{\begin{column}{#1\textwidth} #2 \end{column}}
\newcommand{\colb}[3][.48]{\begin{column}{#1\textwidth} \begin{block}{#2} #3 \end{block} \end{column}}

% Spliting frames in 3 columns
\newcommand{\splitthree}[6]{
  \begin{columns}[T] % align columns
  \begin{column}{#1\textwidth} #4 \end{column} ~~~
  \begin{column}{#2\textwidth} #5 \end{column} ~~~
  \begin{column}{#3\textwidth} #6 \end{column} \end{columns}
}
\newcommand{\frsplitthree}[4][.31]{
  \begin{columns}%[T] % align columns
  \begin{column}{#1\textwidth} #2 \end{column} ~~~
  \begin{column}{#1\textwidth} #3 \end{column} ~~~
  \begin{column}{#1\textwidth} #4 \end{column} \end{columns}
}
\newcommand{\frsplitdiffthree}[7][]{
  \begin{columns}[#1]%[T] % align columns
  \begin{column}{#2\textwidth} #5 \end{column} ~~~
  \begin{column}{#3\textwidth} #6 \end{column} ~~~
  \begin{column}{#4\textwidth} #7 \end{column} \end{columns}
}
\newcommand{\frsplittthree}[4][.32]{
  \begin{columns}[T] % align columns
  \begin{column}{#1\textwidth} #2 \end{column} ~
  \begin{column}{#1\textwidth} #3 \end{column} ~
  \begin{column}{#1\textwidth} #4 \end{column} \end{columns}
}


\newcommand{\typerule}[4][]{\ensuremath{\begin{array}[#1]{c}\textsf{\scriptsize ({#2})} \\#3 \\\hline\raisebox{-3pt}{\ensuremath{#4}}\end{array}}}
\newcommand{\styperule}[3][]{\ensuremath{\begin{array}[#1]{c} #2 \\[0.5mm]\hline\raisebox{-4pt}{\ensuremath{#3}}\end{array}}}
\newcommand{\shrk}{\vspace{-3mm}}

\def\caixa#1{\medskip
  \begin{center}
  \fbox{\begin{minipage}{0.9\textwidth}\protect{#1}\end{minipage}}
  \end{center}}

\newcommand{\mybox}[2][4mm]{
  % \begin{minipage}{#1\textwidth}\begin{block}{}\centering #2\end{block}\end{minipage}}
  \tikz{\node[fill=barcolor!60,align=center,inner sep=#1]{#2};}}
\newcommand{\mycbox}[2][4mm]{
  \begin{center}\mybox[#1]{#2}\end{center}}
  % {\\[-5mm]\centering\mybox[#1]{#2}\\[-5mm]}}

%%%%% Tikz
% \usetikzlibrary{arrows.meta, calc, fit, tikzmark}
\usetikzlibrary{%
  positioning
 ,patterns
 ,arrows
 ,arrows.meta
 ,automata
 ,calc
 ,shapes
 ,fit
 ,tikzmark
 ,fadings
 ,decorations.pathreplacing
 ,plotmarks
% ,pgfplots.groupplots
 ,decorations.markings
 ,shadows
}
% \tikzset{shorten >=1pt,node distance=2cm,on grid,auto,initial text={},inner sep=2pt}
\tikzstyle{aut}=[shorten >=1pt,node distance=2cm,on grid,auto,initial text={},inner sep=2pt]
\tikzstyle{st}=[circle,draw=black,fill=black!10,inner sep=3pt]
\tikzstyle{sst}=[rectangle,draw=none,fill=none,inner sep=3pt]
\tikzstyle{final}=[accepting]

%%% Uppaal-like diagrams
\newcommand{\uppbox}[3][20mm]{\tikz{
  \node[black!15,fill=black!15,minimum width=#1,align=left](title){\textbf{{\footnotesize #2}}};
  \node[black!15,fill=black!15,left,xshift=4mm]at(title.east){\textbf{{\footnotesize #2}}};
  \node[blue!60!cyan,right] at(title.west){\textbf{{\footnotesize #2}}};
  \node[below,inner sep=2mm,fill=white,xshift=2mm](box)at(title.south){\includegraphics[width=#1]{#3}};
  \node[fit=(title)(box),draw=black,inner sep=0pt]{};
}}

\newcommand{\uppboxv}[3][20mm]{\tikz{
  \node[below,inner sep=2mm,fill=white](box){\includegraphics[height=#1]{#3}};
  \coordinate[yshift=5mm](top)at(box.north);
  \node[fit=(top)(box.north west)(box.north east),inner sep=0pt,fill=black!15](title){};
  \node[blue!60!cyan,right] at(title.west){\textbf{{\footnotesize #2}}};
  \node[fit=(title)(box),draw=black,inner sep=0pt]{};
}}

%%% frame for pictures [graphics-options]{content}
\newcommand{\includegraphicsframed}[2][]{\tikz{
  \node[below,inner sep=1mm,fill=barcolor,%xshift=2mm,
    draw=black,
    drop shadow={
      top color=gray,
      bottom color=white,
      %fill=gray,
      opacity=0.2,
      shadow xshift=4pt,
      shadow yshift=-3pt
    }](box){\includegraphics[#1]{#2}};
}}


%% COnfiguring Listings
\lstset{ % basic style
  language=scala,
  basicstyle=\ttfamily\scriptsize,
  breakatwhitespace=true,
  breaklines=true,
  % mathescape,
  % morecomment=[l]{//},
  % morecomment=[n]{/*}{*/},
  % frame=single,                    % adds a frame around the code
  rulecolor=\color{black!40},         % if not set, the frame-color may be changed on line-breaks within not-black text (e.g. comments (green here))
  xleftmargin=1.5mm,
  xrightmargin=1.5mm,
  backgroundcolor=\color{black!5},
  % line numbers
%  numbers=left,  % where to put the line-numbers; possible values are (none, left, right)
 numbersep=5pt, % how far the line-numbers are from the code
 numberstyle=\tiny\color{gray},   
 stepnumber=1,  % the step between two line-numbers. If it is 1 each line will be numbered      
%  xleftmargin=3mm,
%  xrightmargin=1.5mm,
%%%%%
  captionpos=b, % t or b (top or bottom)
  belowcaptionskip=5mm,
%%%%%
  % alsoletter={-},
  % emphstyle=\ttfamily\color{blue}, %\underbar,
  % emphstyle={[2]\ttfamily\color{green!50!black}},
  emphstyle=\bfseries\itshape\color{blue!80!black},       % moreemph={...} - layer keywords
  emphstyle={[2]\itshape\color{red!70!black}},%\underbar} % moreemph={[2]...} - inner keywords
  %
  keywordstyle=\bf\ttfamily\color{red!50!black},
  commentstyle=\sl\ttfamily\color{gray!70},
  % commentstyle=\color{green!60!black},
  stringstyle=\ttfamily\color{green!60!black},
  morestring=[b]",
  % morecomment=[l]{\#},
  frame=single,
  % numberstyle=\tiny, numbers=left, stepnumber=1, firstnumber=1, numberfirstline=true,
  %emph={act,proc,init,sort,map,var,eqn},
  %emph={[2]block,hide,comm,rename,allow,||,<>,sum,&&,=>,true,false},
  literate={\\§}{{{\mbox{\textdollar}}}}1,
  % literate=*{->}{{{\color{red!70!black}$\to$}}}{1}
  %            {.}{{{\color{red!70!black}.}}}{1}
  %            {+}{{{\color{red!70!black}\hspace*{1pt}+\hspace*{1pt}}}}{1}
  %            {|}{{{\color{red!70!black}|}}}{1}
  %            {||}{{{\color{red!70!black}|\!\!|}}}{1}
  %            {*}{{{\color{red!70!black}*}}}{1}
  %            {\#}{{{\color{red!70!black}\#}}}{1}
  %            {&}{{{\color{red!70!black}\&}}}{1}
  %            {:=}{{{\color{red!70!black}:=}}}{1}
  %            {=>}{{{\color{red!70!black}=>}}}{1}
  %            % {>}{{{\color{green!65!black}\hspace*{1pt}>\hspace*{1.5pt}}}}{2}
  %            % {<}{{{\color{green!65!black}\hspace*{1.5pt}<\hspace*{1pt}}}}{2}
  %            % {]}{{{\color{green!65!black}\hspace*{1pt}]\hspace*{1.5pt}}}}{1}
  %            % {[}{{{\color{green!65!black}\hspace*{1.5pt}[\hspace*{1pt}}}}{1}
  %            {>}{{{\color{green!65!black}>}}}{1}
  %            {<}{{{\color{green!65!black}<}}}{1}
  %            {]}{{{\color{green!65!black}]}}}{1}
  %            {[}{{{\color{green!65!black}[}}}{1}
  %  morekeywords={Merger1,Fifo2,Lossy3,Init1,Init2}
  % ,emph={act,proc,init,sort,map,var,eqn}
  % ,emph={[2]block,hide,comm,rename,allow,||,<>,sum,&&,=>}
}

\newcommand{\code}[1]{{\relsize{-1}\ttfamily #1}}
\newcommand{\cod}[1]{{\relsize{-0.5}\color{blue!75!black}\ttfamily #1}}
% \definecolor{mydarkgreen}{rgb}{0,0.6,0}
% \newcommand{\bash}[1]{\lstinline[basicstyle=\ttfamily\relsize{-0.5}\color{mydarkgreen},keywordstyle=\bf\sffamily\color{purple},columns=fullflexible,keepspaces,literate=*]|#1|}

\lstdefinestyle{tiny}{basicstyle=\ttfamily\relsize{-7}}

% Include slides from others
\newcommand{\byothers}[3]{{
\begin{frame}{}~\mycbox{\Large slides by #1\\pages #2}\end{frame}{}
\setbeamercolor{background canvas}{bg=}
\includepdf[pages=#2]{../../others/#3}
}}

\newcounter{mypage}
\newcommand{\fromBook}[6][scale=0.5]{
\setcounter{mypage}{20}
\addtocounter{mypage}{#2}
\begin{tabular}{@{}r@{}}
\includegraphics[page=\themypage, trim = #3 #4 #5 #6, clip, #1]%
  {../../learning-concurrent-scala/learning-concurrent-programming-in-scala.pdf}
  \\[0mm]
  \textcolor{myGray}{$\left[\begin{array}{@{}r@{}}
  ~\\[-7mm]
  \text{{\tiny in \emph{``Learning Concurrent}}}\\[-3mm]
  \text{{\tiny \emph{Programming in Scala''}, pg.\,#2}}
  \\[0mm]
  \end{array}\right]$}
\end{tabular}
}

\newcommand{\fromBookW}[4][width=\textwidth]{
\setcounter{mypage}{20}
\addtocounter{mypage}{#2}
\begin{tabular}{@{}r@{}}
\includegraphics[page=\themypage, trim = 20mm #3 20mm #4, clip, #1]%
  {../../learning-concurrent-scala/learning-concurrent-programming-in-scala.pdf}
  \\[-2mm]
  \textcolor{myGray}{$\left[\begin{array}{@{}r@{}}
  ~\\[-7mm]
  \text{{\tiny in \emph{``Learning Concurrent}
                      \emph{Programming in Scala''}, pg.\,#2}}
  \\[0mm]
  \end{array}\right]$}
\end{tabular}
}

\newcommand{\fromBookWframed}[4][width=\textwidth]{
\setcounter{mypage}{20}
\addtocounter{mypage}{#2}
\begin{tabular}{@{}r@{}}
\includegraphicsframed[page=\themypage, trim = 20mm #3 20mm #4, clip, #1]%
  {../../learning-concurrent-scala/learning-concurrent-programming-in-scala.pdf}
  \\[-2mm]
  \textcolor{myGray}{$\left[\begin{array}{@{}r@{}}
  ~\\[-7mm]
  \text{{\tiny in \emph{``Learning Concurrent}
                      \emph{Programming in Scala''}, pg.\,#2}}
  \\[0mm]
  \end{array}\right]$}
\end{tabular}
}

\newcommand{\fromAuthor}[3][width=\textwidth]{
\begin{tabular}{@{}r@{}}
\includegraphics[#1]%
  {src/img/#3}
  \\[-2mm]
  \textcolor{myGray}{$\left[\begin{array}{@{}r@{}}
  ~\\[-7mm]
  \text{{\tiny in #2}}
  \\[0mm]
  \end{array}\right]$}
\end{tabular}
}


%-------------- template --------------------------------------------------
 %!TEX root=../../1-intro.tex
 \usetheme{metropolis}
\usepackage{appendixnumberbeamer}

%------ Setting lecture info ----------------------------------------------
\newcounter{lectureID}
\stepcounter{lectureID}
\newcommand{\getLecture}{\arabic{lectureID}\xspace}
\newcommand{\setLectureBasic}[1]{
  \title{
    #1
    }
  \author{Nelma Moreira ~~\&~~ \textbf{Jos\'{e} Proen\c{c}a}}
  \institute{CISTER -- U.Porto, Porto, Portugal
            \hfill 
            \begin{tabular}{r@{}}
            \url{https://fm-dcc.github.io/cp2425}
            \end{tabular}
            }
  \date{Concurrent programming (CC3040) 2024/2025}
  % logos of institutions
  \titlegraphic{
    \begin{textblock*}{5cm}(2.0cm,7.00cm)
       \includegraphics[scale=0.14]{src/img/logos/fcup}\hspace*{.85cm}~%
    \end{textblock*}
    \begin{textblock*}{5cm}(6.0cm,7.25cm)
      \includegraphics[scale=0.43]{src/img/logos/cmup}
    \end{textblock*}
    \begin{textblock*}{5cm}(10.4cm,7.45cm)
      \includegraphics[scale=0.16]{src/img/logos/cister}
    \end{textblock*}
  }  
  \frame[plain]{\titlepage}
}
\newcommand{\setLecture}[2]{\setcounter{lectureID}{#1}\setLectureBasic{#1. #2}}

%------ Counters for exercises ----------------------------------------------
\newcounter{cExercise}
\newcommand{\exercise}{\stepcounter{cExercise}Ex.\,\arabic{lectureID}.\arabic{cExercise}:\xspace}
\newcommand{\exerciseBack}{\addtocounter{cExercise}{-1}}
\newcommand{\exerciseAdd}{\stepcounter{cExercise}}
\newcommand{\doExercise}[3][0mm]{\begin{exampleblock}{\exercise #2}\wrap{\rule{0pt}{#1}}#3\end{exampleblock}}
\newcommand{\doSimpleExercise}[2][0mm]{\begin{exampleblock}{}\wrap{\rule{0pt}{#1}}\structure{\textbf{\exercise} #2}\end{exampleblock}}

% Slide
\newenvironment{slide}[1]{\begin{frame}\frametitle{#1}}{\end{frame}}



% Base colors (from metropolis theme)
\definecolor{metDarkBrown}{HTML}{604c38}
\definecolor{metDarkTeal}{HTML}{23373b}
\definecolor{metLightBrown}{HTML}{EB811B}
\definecolor{metLightGreen}{HTML}{14B03D}
\definecolor{metDarkGreen}{HTML}{4f9277}

 

\metroset{numbering=fraction,progressbar=frametitle}

% \setbeamercolor*{structure}{fg=blue!80!black}
\setbeamercolor*{structure}{fg=metLightGreen}

% % \definecolor{MainColour}{rgb}{0., 0.25, 0.8}
% \colorlet{MainColour}{blue!50!black}
% \colorlet{BgColour}{blue!10}
% \colorlet{BarColour}{blue!50!black}

% %\usetheme{CambridgeUS}%{Copenhagen}%{Frankfurt}%{Singapore}%{CambridgeUS}
% \usecolortheme[named=MainColour]{structure} 
% \useoutertheme[subsection=false]{miniframes}
% \useinnertheme{circles}
% %\useinnertheme[shadow=false]{rounded}
% \setbeamertemplate{blocks}[rounded][shadow=false]

% \setbeamercovered{transparent} 
% \setbeamertemplate{navigation symbols}{} %Remove navigation bar
% \setbeamertemplate{footline}[frame number] % add page number
% \setbeamercolor{postit}{fg=MainColour,bg=BgColour}
% \setbeamercolor{structure}{bg=black!10}
% %\setbeamercolor{palette primary}{use=structure,fg=red,bg=green}
% %\setbeamercolor{palette secondary}{use=structure,fg=red!75!black,bg=green}
% \setbeamercolor{palette tertiary}{use=structure,bg=BarColour,fg=white}
% %\setbeamercolor{palette quaternary}{fg=black,bg=green}
% %\setbeamercolor{normal text}{fg=black,bg=white}
% %\setbeamercolor{block title alerted}{fg=red,bg=green}
% %\setbeamercolor{block title example}{bg=black!10,fg=green}
\setbeamercolor{block body}{bg=black!5}

% \setbeamercolor{block title alerted}{bg=red!25}
% \setbeamercolor{block body alerted}{bg=red!10}

% \setbeamercolor{block title example}{bg={rgb:green,2;black,1;white,5}}
\setbeamercolor{block body example}{bg={rgb:green,2;black,1;white,20}}
\setbeamercolor{block body alerted}{bg={metLightBrown!25}}

% \setbeamertemplate{itemize item}{\color{black!10}$\blacksquare$}
\setbeamercolor{itemize item}{fg=metDarkTeal}
\setbeamercolor{itemize subitem}{fg=metDarkTeal}

\setbeamercolor{graybc}{fg=black,bg=black!10}
\newcommand{\myblock}[1]{\begin{beamercolorbox}[dp=1ex,center,rounded=true]%
  {graybc} {\large \textbf{#1}} \end{beamercolorbox}}%

%%%%%%%%%

% \definecolor{barcolor}{rgb}{.65,.79,.92} % FCUP color
\colorlet{barcolor}{metDarkGreen!60}


% Configuring the foot line
\setbeamercolor{author in head/foot}{fg=metDarkTeal, bg=barcolor}%
\setbeamercolor{date in head/foot}{fg=barcolor!75!black}%
\setbeamertemplate{footline}
{
  \leavevmode%
  \hbox{%
  \begin{beamercolorbox}[wd=.4\paperwidth,ht=2.25ex,dp=1ex,center]{author in head/foot}%
    \usebeamerfont{author in head/foot}\insertshortauthor
  \end{beamercolorbox}%
  \begin{beamercolorbox}[wd=.5\paperwidth,ht=2.25ex,dp=1ex,center]{title in head/foot}%
    \usebeamerfont{title in head/foot}\insertsection
  \end{beamercolorbox}%
  \begin{beamercolorbox}[wd=.1\paperwidth,ht=2.25ex,dp=1ex,right]{date in head/foot}%
    \insertframenumber{} / \inserttotalframenumber\hspace*{2ex} 
  \end{beamercolorbox}}%
  \vskip0pt%
}
% No configuration symbols
\setbeamertemplate{navigation symbols}{}


%%%%%%% Custom bar above

\setbeamercolor{frametitle}{bg=barcolor,fg=metDarkTeal}
\setbeamercolor{progress bar}{fg=barcolor!75!black}

%%% Custom bar with a LOGO

\makeatletter
\setbeamertemplate{frametitle}{%
  \nointerlineskip%
  \begin{beamercolorbox}[%
      wd=\paperwidth,%
      sep=0pt,%
      leftskip=\metropolis@frametitle@padding,%
      rightskip=\metropolis@frametitle@padding,%
    ]{frametitle}%
  \metropolis@frametitlestrut@start%
  \insertframetitle%
  \nolinebreak%
  \metropolis@frametitlestrut@end%
  \hfill
  \raisebox{-1.3ex}{\includegraphics[height=4ex,keepaspectratio]{src/img/logos/fcup-bar}}
  \end{beamercolorbox}%
}
\makeatother





%----------------------------------------------------------------------------

\begin{document}

\setLecture{4}{Basic building blocks of concurrency}


\section{Overview}

\begin{frame}{We are here}

  \vspace*{-2mm}

  \begin{block}{Blocks of sequential code running concurrently and sharing memory:}
    
  \begin{itemize}
    \item What is Scala?
    \item Concurrency in Java and its memory model
    \alert{\item Basic concurrency blocks and libraries}
    % \item Futures and promises
    \item \textcolor{gray}{\emph{Futures and Promises}}
    \item \textcolor{gray}{\emph{Data-Parallel Collections}}
    \item \textcolor{gray}{\emph{Reactive Programming (Concurrently)}}
    \item \textcolor{gray}{\emph{Software Transactional Memory}}
    \item Actor model
  \end{itemize}
  \end{block}
\end{frame}



% \fromBookW[scale=0.7]{32}{98mm}{39mm}
% \fromBook[scale=0.65]{32}{43mm}{98mm}{43mm}{39mm}

% \begin{frame}[fragile]\frametitle{Current thread}
% ~\\[-8mm]
% \begin{columns}
% \begin{column}{0.49\textwidth}
% \begin{lstlisting}
% ...
% \end{lstlisting}
% \end{column}
% \begin{column}{0.49\textwidth}
% ...
% \end{column}
% \end{columns}
% \end{frame}


\begin{frame}[t]\frametitle{What we will see}

  \begin{itemize}
    \item Tread pools: Executor and ExecutionContext
    \item Non-blocking synchronisation -- compare-and-set (CAS)
    \item Lazy (concurrent) values
    \item Concurrent collections
    \item Running OS processes
  \end{itemize}


\end{frame}

\section{Existing thread pools in Scala}


\begin{frame}[fragile]\frametitle{Executor interface}
~\\[-8mm]
\begin{columns}
\begin{column}{0.54\textwidth}
\begin{lstlisting}[emph={executor},mathescape]
 Executor executor = $\textit{anExecutor}$;
 executor.execute(new RunnableTask1());
 executor.execute(new RunnableTask2());
 ...
\end{lstlisting}
\begin{lstlisting}[emph={executor}]
import scala.concurrent._
import java.util.concurrent.ForkJoinPool

object ExecutorsCreate extends App {
  val executor = new ForkJoinPool
  executor.execute(new Runnable {
    def run() = log("This task is run asynchronously.")
  })
  Thread.sleep(500) // not needed with fork:=false in SBT
}
\end{lstlisting}
\end{column}
\begin{column}{0.49\textwidth}
\begin{itemize}
  \item \alert{Executor:} can start a new thread, an existing one, or the current one
  \item Abstracts from the management of threads
  \item \structure{ExecutorService:} API that extends Executor with \code{shutdown}
    \begin{itemize}
      \item \code{executor.shutdown} $\to$ executes all tasks and then stops working threads
      \item \code{executor.awaitTermination(...)} $\to$ force termination if, after a given time, the tasks are not completed
    \end{itemize}
\end{itemize}
\end{column}
\end{columns}
\end{frame}


\begin{frame}[fragile]\frametitle{Scala's \texttt{ExecutionContext}}
~\\[-6mm]
\begin{columns}
\begin{column}{0.74\textwidth}
\begin{lstlisting}[emph={execute,sleep,log,global}]
import scala.concurrent._
object ExecutionContextGlobal extends App {
  val ectx = ExecutionContext.global
  ectx.execute(new Runnable {
    def run() = log("Running on the execution context.")
  })
  Thread.sleep(500)
}
\end{lstlisting}
\begin{lstlisting}[emph={execute,sleep,log,global,fromExecutorService}]
object ExecutionContextCreate extends App {
  val pool = new forkjoin.ForkJoinPool(2)
  val ectx = ExecutionContext.fromExecutorService(pool)
  ectx.execute(new Runnable {
    def run() = log("Running on the execution context again.")
  })
  Thread.sleep(500)
}
\end{lstlisting}
\end{column}
\begin{column}{0.36\textwidth}
\begin{itemize}
  \item \code{scala.concurrent}: has \alert{ExecutionContext}
  \item Similar to \structure{Executor} but more Scala specific
  \item often used as implicit parameter
  \item \alert{\code{global}}: default execution context (internally uses a \code{ForkJoinPool})
  \item \code{fromExecutorService}: creates \structure{ExecutionContext} from \structure{ExecutorService}
\end{itemize}
\end{column}
\end{columns}
\end{frame}


\begin{frame}[fragile]\frametitle{Simplifying the execution}
~\\[-8mm]
\begin{columns}[t]
\begin{column}{0.45\textwidth}
Similar to \structure{threads}:

\begin{lstlisting}[emph={thread,run,start}]
def thread(body: =>Unit): Thread = {
  val t = new Thread {
    override def run() = body
  }
  t.start()
  t
}
\end{lstlisting}
\end{column}
\begin{column}{0.56\textwidth}
We now define \alert{\code{execute}}

\begin{lstlisting}[emph={execute,sleep,log,global}]
def execute(body: =>Unit) = ExecutionContext.global.execute(
     new Runnable { def run() = body }
)
// For example:
object ExecutionContextSleep extends App {
  for (i<- 0 until 32) execute {
    Thread.sleep(2000)
    log(s"Task \§i completed.")
  }
  Thread.sleep(10000)
}
\end{lstlisting}
\end{column}
\end{columns}
\end{frame}


\begin{frame}[fragile]\frametitle{Avoid blocking indefinitely}
~\\[-8mm]
\begin{columns}
\begin{column}{0.45\textwidth}
\begin{lstlisting}[emph={execute,sleep,log,global}]
object ExecutionContextSleep
       extends App {
  for (i<- 0 until 32) execute {
    Thread.sleep(2000)
    log(s"Task \§i completed.")
  }
  Thread.sleep(10000)
}
\end{lstlisting}
\end{column}
\begin{column}{0.57\textwidth}
\begin{itemize}
  \item \structure{Expected:} all executions terminate after 2s
  \item \alert{Result:} only some execute after 2s %threads block forever (?!)
  \pause\\[4mm]
  \item Using quad-core CPU with hyper threading
  \item \code{global} has 8 threads in the thread pool
  \pause\\[4mm]
  \item executes tasks in batches of 8
  \item after 2s, 8 tasks print "completed"
  \item after 2s more, 8 more print "completed"
  \item \structure{sleep}: all enter a \structure{timed waiting state}
  \pause\\[4mm]
  \item if T1 \structure{waits} for T10 to \structure{notify}: \alert{blocks indefinitely}
\end{itemize}
\end{column}
\end{columns}
\end{frame}



\section{Lock-free programming}

\begin{frame}[fragile]\frametitle{Avoiding \texttt{syncrhonized} with atomic variables}
~\\[-8mm]
\begin{columns}
\begin{column}{0.57\textwidth}
\begin{itemize}
  \item \alert{atomic variable}: memory location that supports \alert{complex linearizable operations}
  \item ... i.e., \structure{appears to occur atomically}
  \item \code{write} of a volatile operation:
      \\ \structure{simple} linearizable operation
  \item at least two reads and/or writes:
      \\ \structure{complex} linearizable operation
  \pause
  \item \alert{java.util.concurrent.atomic} supports some complex ones:
    \begin{itemize}
       \item AtomicBoolean
       \item AtomicInteger
       \item AtomicLong
       \item AtomicReference  
    \end{itemize}
\end{itemize}
\end{column}
\begin{column}{0.45\textwidth}
~\\
\textbf{Variation of Example 1 (\code{getUniqueId})}
\begin{lstlisting}[emph={execute,sleep,log,global,AtomicLong,incrementAndGet}]
import java.util.concurrent.atomic._

object AtomicUid extends App {
  private val uid =
    new AtomicLong(0L)

  def getUniqueId(): Long =
    uid.incrementAndGet()
  execute {
    log(s"Uid asynchronously: \§{getUniqueId()}")
  }
  log(s"Got a unique id: \§{getUniqueId()}")
}
\end{lstlisting}
\end{column}
\end{columns}
\end{frame}


\begin{frame}[fragile]\frametitle{Compare-And-Set (CAS) -- the \alert{$\heartsuit$} of complex linearizable operations}
~\\[-8mm]
\begin{columns}
\begin{column}{0.48\textwidth}
\begin{itemize}
  \item CAS can be used to implement others:
    \begin{itemize}
      \item getAndSet
      \item decrementAndGet
      \item addAndGet
    \end{itemize}
  \item available in all atomic variables
  \item including AtomicReference[T]
\end{itemize}
\end{column}
\begin{column}{0.54\textwidth}
~\\
\textbf{Long-CAS conceptually equivalent to:}
\begin{lstlisting}[emph={execute,sleep,log,Long,synchronized}]
def compareAndSet(ov: Long, nv: Long):
          Boolean = this.synchronized {
    if (this.get != ov) false else {
      this.set(nv)
      true
  } }
\end{lstlisting}
\textbf{Ref-CAS conceptually equivalent to:}
\begin{lstlisting}[emph={execute,sleep,log,T,synchronized}]
def compareAndSet(ov: T, nv: T):
          Boolean = this.synchronized {
  if (!(this.get eq ov)) false else {
    this.set(nv)
    true
} }
\end{lstlisting}
\end{column}
\end{columns}
\end{frame}


\begin{frame}[fragile]\frametitle{Using CAS}
~\\[-8mm]
\begin{columns}
\begin{column}{0.48\textwidth}
\begin{itemize}
  \item Back to Example~1 (getUniqueId)
  \item Need to keep-on-trying
  \item Looks like busy-waiting, but it is much better
  \item Here: using (cheap) recursion instead of a loop
\end{itemize}
\end{column}
\begin{column}{0.54\textwidth}
~\\
\begin{lstlisting}[emph={execute,sleep,log,compareAndSet,synchronized}]
@tailrec def getUniqueId(): Long = {
  val oldUid = uid.get
  val newUid = oldUid + 1
  if (uid.compareAndSet(oldUid, newUid)) newUid
  else getUniqueId()
}\end{lstlisting}
\end{column}
\end{columns}
\end{frame}


\begin{frame}[fragile]\frametitle{Lock-free programming -- really?}
~\\[-8mm]
\begin{columns}
\begin{column}{0.48\textwidth}
\begin{itemize}
  \item \alert{Lock-free programs}: without locks (with \code{synchronized})
  \item Achieved using \alert{atomic variables} (and some re-trying)
  \item No locks, no deadlocks...
  \pause
  \item (almost):
    \begin{itemize}
      \item lock-free $\Rightarrow$ use atomic variables (for atomicity)
      \item use atomic variables $\cancel\Rightarrow$ lock-free
    \end{itemize}
  \pause
\end{itemize}
\end{column}
\begin{column}{0.54\textwidth}
~\\
\begin{lstlisting}[emph={execute,sleep,log,compareAndSet,mySynchronized,AtomicBoolean}]
object AtomicLock extends App {
  private val lock = new AtomicBoolean(false)
  def mySynchronized(body: =>Unit): Unit = {
    while (!lock.compareAndSet(false, true)) {}
    try body finally lock.set(false)
  }
  var count = 0
  for (i<- 0 until 10) execute { mySynchronized { count += 1 } }
  Thread.sleep(1000)
  log(s"Count is: \§count")
}
\end{lstlisting}
\end{column}
\end{columns}
\end{frame}  


\begin{frame}\frametitle{Lock-freedom definition}

\begin{alertblock}{Lock-freedom}
  Given a set of threads and an operation \alert{OP}.
  \\
  \alert{OP} is \structure{lock-free} if at least one thread always completes \alert{OP} after a finite number of steps, regardless of the speed at which different threads progress.
\end{alertblock}
\end{frame}


\begin{frame}\frametitle{One more example: Concurrent filesystem}
    \begin{itemize}
      \item \textbf{Example 1:} getUniqueId()
      \item \textbf{Example 2:} Logging Bank Transfers
      \item \textbf{Example 3:} Thread pool
      \item \textbf{Example 4:} Batman
      \item \alert{\textbf{Example 5:} Concurrent filesystem}
    \end{itemize}
\end{frame}


\begin{frame}\frametitle{Concurrent filesystem}
  \splittwo{0.35}{0.68}{
    \begin{block}{Filesystem API}
      \alert{T1} is creating \structure{F}:
      \\\alert{T2} cannot copy or delete \structure{F}
      \\[4mm]\alert{T1}\,\&\,\alert{T2} are copying \structure{F}:
      \\\alert{T3} cannot delete \structure{F}
      \\[4mm]\alert{T1} is deleting \structure{F}:
      \\\alert{T2} cannot copy nor delete \structure{F}
    \end{block}
  }{
    \fromBook[scale=0.8]{79}{40mm}{121mm}{40mm}{67mm}
  }
\end{frame}


\begin{frame}[fragile]\frametitle{Concurrent filesystem -- Scala data types}
~\\[-8mm]
\begin{columns}
\begin{column}{0.35\textwidth}
    \begin{block}{Filesystem API}
      \alert{T1} is creating \structure{F}:
      \\\alert{T2} cannot copy or delete \structure{F}
      \\[4mm]\alert{T1}\,\&\,\alert{T2} are copying \structure{F}:
      \\\alert{T3} cannot delete \structure{F}
      \\[4mm]\alert{T1} is deleting \structure{F}:
      \\\alert{T2} cannot copy nor delete \structure{F}
    \end{block}
\end{column}
\begin{column}{0.68\textwidth}
\begin{lstlisting}[emph={execute,sleep,log,compareAndSet,State,AtomicReference}]
class Entry(val isDir: Boolean) {
  val state = new AtomicReference[State](new Idle)
}

sealed trait State
class Idle extends State
class Creating extends State
class Copying(val n: Int) extends State
class Deleting extends State
\end{lstlisting}
\end{column}
\end{columns}
\end{frame}


\begin{frame}[fragile]\frametitle{Deleting and Copying}
\alert{Deleting:} \textbf{\structure{\underline{prepare}}} (checks for permission) then \structure{delete} (perform delete)

\alert{Copying:} \structure{aquire} (get permission); \structure{copy} (perform action); then \structure{release} (give permission)
\end{frame}


\begin{frame}[fragile]\frametitle{Prepare for deleting}

\begin{lstlisting}[emph={execute,sleep,log,logMessage,compareAndSet,State,AtomicReference}]
@tailrec
private def prepareForDelete(entry: Entry): Boolean = {
  val  s0 = entry.state.get
  s0 match {
    case i: Idle =>
      if (entry.state.compareAndSet(s0, new Deleting)) true
      else prepareForDelete(entry)
    case c: Creating =>
      logMessage("File currently created, cannot delete."); false
    case c: Copying =>
      logMessage("File currently copied, cannot delete."); false
    case d: Deleting =>
      false
  }
}
\end{lstlisting}

\cod{logMessage}: presented later -- similar to \cod{log}, but stores the log message
\end{frame}


\begin{frame}\frametitle{Bad copy -- the ABA problem}
\alert{``A\structure{B}A'' problem}: two readings of the same value \alert{A} lead to believe that \structure{B} was never present (type of race condition)

\bigskip

Illustrated by a bad acquire-release for \cod{Copying}, using a mutable \cod{n} in:
\\~~\alert{\texttt{Copying(\structure{var} n: Int)}}
\end{frame}


\begin{frame}[fragile]\frametitle{Bad code -- acquire/release Copying}
\label{slide:badcopy}
    
\begin{lstlisting}[emph={execute,sleep,log,logMessage,compareAndSet,releaseCopy}]
def releaseCopy(e: Entry): Copying = e.state.get match {
  case c: Copying =>
    val nstate = if (c.n == 1) new Idle else new Copying(c.n - 1)
    if (e.state.compareAndSet(c, nstate)) c
    else releaseCopy(e)
}
\end{lstlisting}
\begin{lstlisting}[emph={execute,sleep,log,logMessage,compareAndSet,acquireCopy}]
def acquireCopy(e: Entry, c: Copying) = e.state.get match {
  case i: Idle =>
    c.n = 1
    if (!e.state.compareAndSet(i, c)) acquireCopy(e, c)
  case oc: Copying =>
    c.n = oc.n + 1
    if (!e.state.compareAndSet(oc, c)) acquireCopy(e, c)
}
\end{lstlisting}
Optimization: reusing previous \cod{Copying} if possible
\end{frame}


% \begin{frame}\frametitle{Bad trace -- T1\&T2 release; T3\&T1 aquire -- (T2 is slow)}
\begin{frame}\frametitle{Bad trace -- T1 release (T2 starts rel.); T3 acquires; T1 acquire; T2 releases wrongly}
    \fromBookW{82}{151mm}{33mm}
\end{frame}


\begin{frame}\frametitle{Some guidelines to avoid the ABA problem}
  \begin{itemize}
    \item use fresh objects in \cod{AtomicReference}
    \item use immutable objects in \cod{AtomicReference}
    \item avoid re-assigning the same value to an atomic variable
    \item only increment or decrement values of numeric atomic variables (if possible)
  \end{itemize}


\end{frame}


\section{Lazy values}

\begin{frame}[fragile]\frametitle{Lazy values can cause deadlocks}
~\\[-8mm]
\begin{columns}
\begin{column}{0.48\textwidth}
\begin{itemize}
  \item \alert{\textbf{lazy values:}} initialized when read for the first time
  \item these should not depend-on/modify state (non-determinism)
  \item code in \structure{singleton objects}: lazy execution
  \item under the hood: first write uses a lock -- to ensure at most a thread initialises a lazy value
  \item \structure{stack overflow} (sequential code)\\~~can become\\
        \alert{deadlock} (concurrent code)

\end{itemize}
\end{column}
\begin{column}{0.54\textwidth}
~\\
\begin{lstlisting}[emph={execute,sleep,log,compareAndSet,mySynchronized,AtomicBoolean}]
object LazyValsCreate extends App {
  var x = 5
  lazy val y = x+2
  execute {log(s"Wrk: y = $y")}
  x = 10
  log(s"Main: y = $y")
  // y = 7 or 12 in both cases
  Thread.sleep(500)
}
\end{lstlisting}
\begin{lstlisting}[emph={execute,sleep,log,compareAndSet,mySynchronized,AtomicBoolean}]
object LazyValsDeadlock extends App {
  object A { lazy val x: Int = B.y }
  object B { lazy val y: Int = A.x }
  execute { B.y }
  A.x
}
\end{lstlisting}
\end{column}
\end{columns}
\end{frame}

\section{Concurrent (mutable) collections}


\begin{frame}[fragile]\frametitle{Default mutable collections are not concurrent}


~\\[-8mm]
\begin{columns}
\begin{column}{0.48\textwidth}
\begin{itemize}
  \item Naive code: arbitrarily returns different results and exceptions
  \item Corruption of the internal state
  \item Possible fixes:
    \begin{itemize}
      \item immutable collections + atomic variables
      \\\only<2>{\structure{(does not scale)}}
      \item mutable collections + synchronized
      \\\only<2>{\structure{(assuming collections do not block; may not scale)}}
      \item dedicated libraries
      \\\only<2>{\structure{(far better performance and scalability)}}
    \end{itemize}
\end{itemize}
\end{column}
\begin{column}{0.54\textwidth}
~\\
\begin{lstlisting}[emph={execute,sleep,log,compareAndSet,asyncAdd}]
import scala.collection._
object CollectionsBad extends App {
  val buffer = mutable.ArrayBuffer[Int]()
  def asyncAdd(numbers: Seq[Int]) = execute {
    buffer ++= numbers
    log(s"buffer = $buffer")
  }
  asyncAdd(0 until 10)
  asyncAdd(10 until 20)
  Thread.sleep(500)
}
\end{lstlisting}
\end{column}
\end{columns}
\end{frame}


\begin{frame}[t]\frametitle{Some concurrent collections}
    
  \begin{itemize}
    \item Concurrent queues 
      \begin{itemize}
         \item java.util.concurrent.BlockingQueue \structure{interface}
        \item ...ArrayBlockingQueue \alert{class} (bounded)
        \item ...LinkedBlockingQueue \alert{class} (unbounded)
       \end{itemize} 
    \item Concurrent Sets and Maps
      \begin{itemize}
        \item scala.collection.concurrent.Map \structure{trait}
        \item java.util.concurrent.ConcurrentHashMap \alert{class}
      \end{itemize}
  \end{itemize}
\end{frame}

\begin{frame}[t]\frametitle{BlockingQueue API}
    \centering
    \fromBook[scale=1.2]{90}{20mm}{24mm}{59mm}{172mm}
\end{frame}

\begin{frame}[fragile]\frametitle{Back to Example 5: logging in our concurrent filesystem}
    
We will compile a queue of messages when \alert{logging} messages in our file system


\begin{lstlisting}[emph={execute,sleep,log,compareAndSet,logger,messages,logMessage}]
class FileSystem(...) {
  ...
  private val messages = new LinkedBlockingQueue[String]
  val logger = new Thread {
    setDaemon(true)
    override def run() = while (true) log(messages.take())
  }
  logger.start()
  def logMessage(msg: String): Unit = messages.offer(msg)
}
\end{lstlisting}

\begin{lstlisting}[emph={execute,sleep,log,compareAndSet,logMessage}]
...
val fileSystem = new FileSystem(".") // to be defined later
fileSystem.logMessage("Testing log!")
\end{lstlisting}
\end{frame}


\begin{frame}\frametitle{Note on iterators}

\begin{itemize}
  \item concurrentQueue.iterator
  \item can produce inconsistent results
  \item while traversing and modifying, the iterator can be updated
  \item (heavier) exceptions create a copy when producing an iterator
\end{itemize}
\end{frame}


\begin{frame}[fragile]\frametitle{Example 5: files as a Map in our FileSystem}

\begin{lstlisting}[emph={execute,sleep,log,ConcurrentHashMap,files,Entry,iterateFiles}]
import scala.collection.convert.decorateAsScala._
import java.io.File
import org.apache.commons.io.FileUtils // needs "commons-io" in build.sbt

class FileSystem(val root: String) {
  val rootDir = new File(root)
  val files: concurrent.Map[String, Entry] =
    new ConcurrentHashMap().asScala
  for (f <- FileUtils.iterateFiles(rootDir, null, false).asScala)
    files.put(f.getName, new Entry(false))

  ...
}
\end{lstlisting}
\end{frame}


\begin{frame}[fragile]\frametitle{Deleting a file}

Recall the \code{prepareForDelete(entry)}

\begin{lstlisting}[emph={execute,sleep,log,deleteFile,files,prepareForDelete,deleteQuietly}]
...
def deleteFile(filename: String): Unit = {
  files.get(filename) match {
    case None =>
      logMessage(s"Path '$filename' does not exist!")
    case Some(entry) if entry.isDir =>
      logMessage(s"Path '$filename' is a directory!")
    case Some(entry) => execute {
      if (prepareForDelete(entry))
        if (FileUtils.deleteQuietly(new File(filename)))
          files.remove(filename)
    }
  }
}
\end{lstlisting}
\end{frame}


\begin{frame}\frametitle{Some complex linearizable methods of concurrent Map}
    \centering
    \fromBookW[scale=0.8]{95}{34mm}{142mm}

    \splittwo{0.48}{0.48}{
      \begin{itemize}
        \item These use ``equals'' instead of the reference (which CAS does)
        \item Avoid \code{null} as key or valye (often used as special values)
      \end{itemize}
    }{
      \begin{itemize}
        \item Methods \code{+=}, \code{-=}, \code{put}, \code{update}, \code{get}, \code{apply}, \code{remove} are (non-complex) linearizable
      \end{itemize}
    }
    % (Using "equals" instead of the reference, as with CAS)
\end{frame}


\section*{Wrapping up our Filesystem (Example 5)}

\begin{frame}[fragile]\frametitle{Copying in our Filesystem}
    Recall our \alert{broken} \structure{aquireCopy}/\structure{releaseCopy} methods (\alert{ABA problem}) -- slide\ref{slide:badcopy}

\begin{lstlisting}[emph={execute,sleep,log,Copying,aquire,compareAndSet}]
@tailrec
private def acquire(entry: Entry): Boolean = {
  val s0 = entry.state.get
  s0 match {
    case _: Creating | _: Deleting =>
      logMessage("File inaccessible, cannot copy."); false
    case i: Idle =>
      if (entry.state.compareAndSet(s0, new Copying(1))) true
      else acquire(entry)
    case c: Copying =>
      if (entry.state.compareAndSet(s0, new Copying(c.n+1))) true
      else acquire(entry)
  }
}
\end{lstlisting}
\end{frame}


\begin{frame}[fragile]\frametitle{Copying in our Filesystem}
    Same CAS retry-approach for releasing.

\begin{lstlisting}[emph={execute,sleep,log,Copying,release,compareAndSet}]
@tailrec
private def release(entry: Entry): Unit = {
  val s0 = entry.state.get
  s0 match {
    case c: Creating =>
      if (!entry.state.compareAndSet(s0, new Idle)) release(entry)
    case c: Copying =>
      val nstate = if (c.n == 1) new Idle else new Copying(c.n-1)
      if (!entry.state.compareAndSet(s0, nstate)) release(entry)
  }
}
\end{lstlisting}
\end{frame}


\begin{frame}[fragile]\frametitle{Copying in our Filesystem}
    Finally: wrapper for copying a file.

\begin{lstlisting}[emph={execute,sleep,log,Copying,copyFile,acquire, release}]
def copyFile(src: String, dest: String): Unit = {
  files.get(src) match {
    case Some(srcEntry) if !srcEntry.isDir => execute {
      if (acquire(srcEntry)) try {
        val destEntry = new Entry(isDir = false)
        destEntry.state.set(new Creating)
        if (files.putIfAbsent(dest, destEntry) == None) try {
          FileUtils.copyFile(new File(src), new File(dest))
        } finally release(destEntry)
      } finally release(srcEntry)
    }
} }
\end{lstlisting}
\end{frame}



\section{Creating and handling processes}

\begin{frame}\frametitle{Beyond JVM}

\begin{itemize}
  \item \structure{So far:} run in a single JVM
  \item \structure{Now:} run processes outside JVM
  \item \structure{Why:}
  \pause
    \begin{itemize}
      \item Some programs do not exist in Scala/Java
      \item \pause Want to sandbox untrusted code
      \item \pause Performance (running independent code)
    \end{itemize}
  \pause
  \item Using the \alert{\code{scala.sys.process}} package
\end{itemize}
\end{frame}

\begin{frame}[fragile]\frametitle{Using processes - examples}
    
\begin{lstlisting}[emph={execute,sleep,log,destroy,!}]
import scala.sys.process._
object ProcessRun extends App {
  val command = "ls"
  val exitcode = command.! // run process (with side effects)
  log(s"command exited with status $exitcode") }
\end{lstlisting}

\begin{lstlisting}[emph={execute,sleep,log,destroy,!}]
def lineCount(filename: String): Int = {
  val output = s"wc $filename".!! // run and retreive stdout
  output.trim.split(" ").head.toInt }
\end{lstlisting}

\begin{lstlisting}[emph={execute,sleep,log,destroy}]
object ProcessAsync extends App {
  val lsProcess = "ls -R /".run() // run and returns a Process object
  Thread.sleep(1000)
  log("Timeout - killing ls!")
  lsProcess.destroy() } // kill a slow process
\end{lstlisting}

{\small\url{https://www.scala-lang.org/api/2.13.x/scala/sys/process/ProcessBuilder.html}}

\end{frame}


\begin{frame}\frametitle{Wrapping up ``concurrency blocks''} %{Summary of concurrency blocks covered}
  \splittwo{0.56}{0.44}{
    \begin{itemize}
      \item \texttt{\color{blue!75!black}executor.execute(...)}
      \item lock-free programming with atomic variables
      \item \texttt{\color{blue!75!black}av.compareAndSet(...)}
      \item ABA problem
      \item Lazy values \& ``lazy'' objects
      \item \texttt{\color{blue!75!black}java.util.concurrent.BlockingQueue}
      \item \texttt{\color{blue!75!black}scala.collection.concurrent.Map}
      \item \textcolor{gray}{\emph{weakly consistent iterators}}
      \item \textcolor{gray}{\emph{custom concurrent data structures}}
    \end{itemize}
  }{
    \begin{itemize}
      \item Filesystem example
      \item Processes outside JVM
    \end{itemize}
    \pause
    \begin{alertblock}{Next}
      \begin{itemize}
        \item \textcolor{gray}{\emph{Futures and Promises}}
        \item \textcolor{gray}{\emph{Data-Parallel Collections}}
        \item \textcolor{gray}{\emph{Reactive Programming (Concurrently)}}
        \item \textcolor{gray}{\emph{Software Transactional Memory}}
        \item \alert{Actors}
      \end{itemize}
    \end{alertblock}
  }
\end{frame}


\end{document}
